% CE7454 Project 2, FY2026. Working draft updated 10 November 2026 (SGT).
% FY2026 handout: four-page CVPR report; no explicit appendix/reference exemption found.
% Keep four total pages pending instructor clarification. No appendix is assumed exempt.
% Source: reporting fix 18fb9b2; shared_hawker_20261108_k4 at f025ee0. Independent evaluation pending.
% Previous live proposal is preserved in Overleaf History with a dated label.
\documentclass[10pt,twocolumn,letterpaper]{article}
\usepackage[pagenumbers]{cvpr}
%% This file contains a number of tweaks that are typically applied to the main document.
%% They are not enabled by default, but can be enabled by uncommenting the relevant lines.

%%
%% Inline annotations; for predefined colors, refer to "dvipsnames" in the xcolor package:
%% https://tinyurl.com/overleaf-colors
%%
\newcommand{\red}[1]{{\color{red}#1}}
\newcommand{\todo}[1]{{\color{red}#1}}
\newcommand{\TODO}[1]{\textbf{\color{red}[TODO: #1]}}
%%
%% disable for camera ready / submission by uncommenting these lines  
%%
% \renewcommand{\TODO}[1]{}
% \renewcommand{\todo}[1]{#1}

%%
%% work harder in optimizing text layout. Typically shrinks text by 1/6 of page, enable
%% it at the very end of the writing process, when you are just above the page limit
%%
% \usepackage{microtype}

%%
%% fine-tune paragraph spacing
%%
% \renewcommand{\paragraph}[1]{\vspace{.5em}\noindent\textbf{#1.}}

%%
%% globally adjusts space between figure and caption
%%
% \setlength{\abovecaptionskip}{.5em}


%%
%% Allows "the use of \paper to refer to the project name"
%% with automatic management of space at the end of the word
%%
% \usepackage{xspace}
% \newcommand{\paper}{ProjectName\xspace}

%%
%% Commonly used math definitions
%%
% \DeclareMathOperator*{\argmin}{arg\,min}
% \DeclareMathOperator*{\argmax}{arg\,max}

%%
%% Tigthen underline
%%
% \usepackage{soul}
% \setuldepth{foobar}
\usepackage{tikz}
\usetikzlibrary{arrows.meta}
\definecolor{cvprblue}{rgb}{0.21,0.49,0.74}
\usepackage[breaklinks,colorlinks,allcolors=cvprblue]{hyperref}
\def\paperID{}
\def\confName{CVPR}
\def\confYear{2026}
\newcommand{\draftnote}[1]{\par\noindent{\color{blue}\small\textbf{Draft:} #1}\par}
\title{SgFood-Edit: Check-and-Retry Meal Previews with a Vision-Language Checker}
% GitHub-safe snapshot: author contact/student identifiers remain in private Overleaf.
\author{CE7454 Project 2 Group 12\\Nanyang Technological University\\
\textit{Author details maintained in the private submission draft}}
\begin{document}
\maketitle
\begin{abstract}
SgFood-Edit couples a frozen BrushNet-SDXL editor with a Qwen2.5-VL checker for failure-directed meal-photo editing. On 20 retrospectively shared photographs (34 tasks per prompt condition), revised-prompt retry passes the selecting checker on 28/34 outputs versus 26/34 for best-of-four. The task effects have opposite signs; the exploratory 95\% photo-cluster interval for the overall difference is $[-5.7,+17.1]$ percentage points. Fixed-prompt early stopping also passes 26/34 with the same 68 aggregate generations as retry, so fewer generations alone do not demonstrate useful feedback. Supplemental CLIP selection passes 22/34; unchanged photos (0/34) and a fixed recolouring control (0/14) are rejected. These diagnostics do not establish independent visual quality, nutritional correctness or health benefit.
\end{abstract}

\section{Introduction}
Singapore's National Nutrition Survey 2025 reports refined-staple carbohydrate intake rising from 154 to 166\,g/day in 2019--2025~\cite{hpbNNS2025}, measured as nutrient, not food weight. HealthHub recommends gradual wholegrain substitution~\cite{healthhubWholegrains}. We study faithful meal-change previews, not effects on awareness or food choices.

Plausible edits can replace the staple, merely recolour rice or change nearby food. Input similarity and outside-mask pixel equality cannot establish success: an unchanged image performs no edit.

A vision-language model (VLM) checks each frozen-editor candidate and directs the next attempt by its main failure. Three questions guide the study: \textbf{Q1}, does targeted retry outperform selecting the best of $K$ candidates at the same maximum generation budget? \textbf{Q2}, how much benefit comes from correcting the prompt rather than the retry loop? \textbf{Q3}, how often does a checker favour outputs that an independent evaluator rejects? Our auditable controller separates prompt revision, selection, stopping and adaptive retry. Q1--Q2 are addressed only through checker diagnostics; independent quality and Q3 remain open.

\section{Related Work}
\paragraph{Controlled image editing.}
Latent diffusion~\cite{rombach2022latent} and SDXL~\cite{podell2024sdxl} provide the backbone; BrushNet~\cite{ju2024brushnet} adds masked-image features for inpainting. Prior work studies caloric targets~\cite{yamamoto2025calorie} and food-advertising backgrounds~\cite{feng2025forage}. We instead change one meal component while retaining the dish, without inferring calories.

\paragraph{Feedback and evaluation.}
Self-Refine~\cite{madaan2023selfrefine} studies feedback for language-model outputs. Idea2Img~\cite{yang2024idea2img} and SLD~\cite{wu2024sld} use multimodal or language-model feedback to revise text-to-image generation iteratively. InstructPix2Pix~\cite{brooks2023instructpix2pix} follows image-editing instructions; MagicBrush~\cite{zhang2023magicbrush} supplies manually annotated instruction-edit pairs. Ma et al.~\cite{ma2025scaling} scale diffusion inference through verifier-guided noise search, establishing selection and search as relevant precedents. We test whether failure-conditioned updates add value beyond sampling and stopping for food edits. VIEScore~\cite{ku2024viescore} reports evaluation difficulties on editing tasks, motivating external assessment of Qwen2.5-VL-selected outputs~\cite{bai2025qwen25vl}. CLIP~\cite{radford2021clip} is an alternative selector, not ground truth.

\section{Method}
\subsection{Task and Frozen Editing Stack}
The input is a meal photograph $x$, a target mask $m$ over its staple, and one of two requests: \textbf{T1-B}, add broccoli while retaining the original rice or noodles; \textbf{T2}, depict brown rice in place of white rice. The implementation uses SDXL base 1.0 with the BrushNet-SDXL segmentation-mask checkpoint and a local Qwen2.5-VL-7B-Instruct checker loaded in 4-bit precision. No model weights are trained. The editor and checker run on separate V100 GPUs.

BrushNet operates on a crop around the editable region at $1024\times1024$ resolution. A feathered mask $\tilde m$ pastes the result back:
\begin{equation}
 y=\tilde m\odot g(x,m,t)+(1-\tilde m)\odot x.
\end{equation}
Source pixels are unchanged wherever $\tilde m=0$; an outside-only band therefore cannot assess seam quality.



\subsection{Checker and Retry Policy}
The checker asks fixed multiple-choice questions about edit completion, preservation, realism and unrequested additions. Scores are expectations over answer-token probabilities, not calibrated probabilities of success. Let $e,p\in[0,2]$ denote edit and preservation scores, $r\in[1,5]$ realism, and $a\in[0,1]$ unrequested additions. The audited policy ranks candidates by
\begin{equation}
 S=e/2+p/4+(r-1)/8-a/2.
\end{equation}
Its pass rule requires $e,p\geq1.5$, $r\geq3.0$ and $a\leq0.5$. The existing T2 rubric requires complete visual replacement of white rice; a mixed result receives partial credit. This operational target is distinct from HealthHub's gradual half-portion substitution and does not certify wholegrain content.

The initial candidate uses 20 diffusion steps and guidance 7.5. If edit completion fails, the policy advances through a task-specific positive/negative prompt ladder; after the ladder is exhausted, guidance can increase to 10.0 and 12.5. Unrequested additions, realism or preservation failures trigger another seed at base guidance. Failure priority is edit, additions, realism, then preservation. Generation stops when the last candidate passes or budget $K$ is exhausted. The implementation returns the first candidate attaining the highest $S$ among those tried, which need not be the candidate that triggered stopping.

\subsection{Comparisons and Prompt Control}
Single-shot returns the first candidate. Best-of-$K$ checker selects the highest $S$ among $K$ base-setting candidates. Adaptive retry uses failure-dependent updates and stopping. Fixed-prompt seed-only stopping replays baseline candidates (Sec.~\ref{sec:budget}); supplemental CLIP and recolouring controls appear in Sec.~\ref{sec:supplemental}.

The original brown-rice prompt used negative instructions within the positive prompt, and a retry extension exceeded the editor's 77-token limit. The revised prompts use shorter positive descriptions and separate negative prompts. The completed policies are evaluated under both prompt conditions to separate prompt and policy effects. Prompt selection, checker thresholds and retry rules use development images only. On the 35 development tasks, first-image and selected-retry passes rise from 12 and 18 with the original prompts to 24 and 30 with the revised prompts.

\section{Experiments}
\subsection{Data and Reproducibility}
Development uses 20 meal photos (15 rice, five noodle) and 35 tasks. Baseline assets and development photos are reused from PlateEdit, a prior CE6190 course project, whose baseline includes FoodSG~\cite{zheng2023foodsg}. CE7454 contributes inference-time control and evaluation, not a new generator.

The separate \emph{shared-photo retrospective robustness} sample contains 20 author-owned photos (14 rice, six noodle), giving 34 cases per prompt condition and 68 condition--case runs. It reuses the CE6190 S5 collection and prior single-author mask QA; it is not an independent held-out test. No byte/normalised-pixel matches or perceptual-hash distances $\leq8$ were found against the 40-photo FoodSG baseline; this does not establish independence. All masks were retained unchanged: 13 accepted and seven flagged. Mask quality is confounded with dish type: the accepted group has 12 rice and one noodle photo; the flagged group has two rice and five noodle photos.

Both prompt conditions were locked before inference on 8 November ($K=4$). Hashes of all inputs, model files and the 365 candidates were audited, and selection and stopping were recomputed without mismatches.
% Submission checklist in README: confirm authors/student IDs and cross-course reuse permission.

\subsection{Evaluation Protocol and Its Limits}
No real ratings have been collected. Planned revised-prompt retry/best-of-four ratings use matched photos, masks and first seeds. The staged analyzer defines success per rater: $e=p=2$, $r\geq3$ (1--5), $a=0$, then averages binary outcomes within cases. Roster and exposure records must be fixed before collection.

Planned intervals resample photos with all their tasks and ratings; dimension scores and rater agreement are reported. False acceptance among checker passes differs from false-positive rate among author failures. Missing/Unsure stay unknown, with bounds; preferences, ties and neither stay separate. Authors are non-naive; another VLM is supplementary.



\subsection{Shared-Photo Checker Diagnostics}
Table~\ref{tab:shared} reports the selected outputs' own checker verdicts. These are descriptive selector diagnostics, not human success rates. With revised prompts, retry exceeds best-of-four by two cases overall (5.9 percentage points), but the task pattern differs: T1-B gains three, whereas T2 loses one. Paired outcomes are 25 passes under both methods, three retry-only passes, one best-of-four-only pass and five failures under both. An exploratory photo-cluster bootstrap gives a 95\% interval of $[-5.7,+17.1]$ points (10,000 resamples; seed 20261108), retaining each photo's tasks and recomputing the pooled rate. The interval includes zero. Case-level exact McNemar gives $p=0.625$ but ignores within-photo dependence. These same-checker diagnostics establish neither external superiority nor freedom from selection bias.

\begin{table}[t]
\centering\small
\setlength{\tabcolsep}{3.5pt}
\begin{tabular}{@{}llrrrr@{}}
\toprule
Prompt & Task & Single & CLIP & Best-of-4 & Retry\\
\midrule
Original & T1-B & 10/20 & -- & 14/20 & 17/20\\
 & T2 & 0/14 & -- & 1/14 & 3/14\\
 & Total & 10/34 & -- & 15/34 & 20/34\\
\midrule
Revised & T1-B & 10/20 & 12/20 & 14/20 & 17/20\\
 & T2 & 10/14 & 10/14 & 12/14 & 11/14\\
 & Total & 20/34 & 22/34 & 26/34 & 28/34\\
\bottomrule
\end{tabular}
\caption{Checker passes on 20 shared photos, 34 tasks per condition, $K=4$. CLIP is supplemental (revised only). Best-of-four and retry select with the evaluating checker; counts are not independent quality estimates. All flagged masks remain included.}
\label{tab:shared}
\end{table}

\begin{figure}[t]
\centering
\IfFileExists{qualitative_examples.pdf}{
\includegraphics[width=\linewidth]{qualitative_examples.pdf}
}{\fbox{\parbox[c][0.27\textheight][c]{0.92\linewidth}{\centering
Private qualitative figure omitted from GitHub.\\
See the reviewed report for the actual image panels.}}}
\caption{Revised-prompt checker-selected outputs, without human success labels. Before visual inspection, the first lexicographic photo ID with non-identical outputs was selected for each task. Full frozen images are shown; these examples do not estimate success rates.}
\label{fig:qualitative}
\end{figure}

Figure~\ref{fig:qualitative} shows actual selected outputs. In the first row, noodle/soup appearance changes alongside the broccoli addition; in the second, changed rice appearance cannot certify wholegrain content.


Prompt revision changes the T2 single-shot count from 0/14 to 10/14, while T1-B counts are unchanged. Comparing original single-shot directly with revised retry would therefore conflate prompt and policy effects. Within revised prompts, retry adds eight checker-passing outputs over single-shot but only two over best-of-four. External ratings are needed to establish visual benefit and audit checker errors. Error on selected outputs alone does not isolate inflation caused by selection.

\subsection{Budget and Early-Stopping Controls}\label{sec:budget}
Budgets $K\in\{1,2,4\}$ are replayed from archived sequences with the original first-maximum rule, not independent runs. Retry averages 1.00, 1.41 and 2.00 generations, respectively; best-of-$K$ always uses $K$. At $K=2$, both pass 23/34; a net checker difference appears only at $K=4$. These generation counts exclude scoring overhead and do not establish equal per-case cost or latency.



A stronger stopping control replays the fixed-prompt baseline seeds, stopping at the first checker pass or $K=4$, and returns the highest-$S$ candidate seen. It passes 26/34 revised-prompt cases using 68 generations in total, versus 28/34 and 68 for adaptive retry. Thus early stopping already explains the mean reduction from four to two generations; adaptive feedback's additional visual benefit remains unverified. With original prompts, seed-only stopping gives 15/34 passes and 99 generations, versus adaptive retry's 20/34 and 93.

Five original and seven revised retry/best-of-four pairs are byte-identical and cannot favour either output.


\subsection{Supplemental Selector and Stress Controls}\label{sec:supplemental}
These retrospective controls were specified after the original study and locked before execution on 10 November. CLIP ViT-L/14 ranks all 136 revised base candidates by full-image target-text cosine, using official preprocessing, FP32 and first-maximum seed-order ties, not a target-minus-source margin. Fixed targets were ``A photograph of a Singapore hawker meal with clearly visible cooked broccoli added to the rice or noodles.'' and ``A photograph of a Singapore hawker meal with light brown wholegrain rice.'' It passes 22/34 under Qwen (Table~\ref{tab:shared}), versus 26/34 for checker selection: $-11.8$ points, exploratory photo-cluster 95\% interval $[-22.2,-2.9]$, using the same bootstrap. Selection optimises the same checker, so this is not independent visual-quality evidence.

Unedited sources pass 0/34 (T1-B: 0/20; T2: 0/14). A separate non-generative control fixes HSV hue to $30^\circ$, saturation to 0.35 and value to $0.82V$ inside each of the 14 original rice masks, with half-up RGB rounding and no feathering or per-image tuning. Outside-mask pixels are unchanged. Joint passes are 0/14; edit, preservation, realism and extras criteria pass 0, 14, 1 and 3 cases. Several controls look conspicuously flat. Rejection of this one transform does not establish general recolouring detection or a human false-acceptance rate. These controls add no diffusion calls to the original 365 candidates.

\subsection{Remaining Evidence}
Per-output ratings of the 34 revised pairs remain essential. Fixed-prompt stopping needs one additional output; CLIP needs 23 outside those pairs. Ratings transfer only for identical pixels, source and task, never pair preferences. Human-judged recovery/reversal examples and new-photo confirmation remain open; mask-stratum comparisons are descriptive.

\section{Limitations and Responsible Use}
The small retrospective sample cannot establish generalisation. Mask strata are confounded with food type, authors may be non-naive, and another VLM may share checker biases. Appearance cannot verify wholegrains, a balanced meal, sodium, calories or portion mass. Interfaces should show originals, label generated previews and disclose failures. Public understanding and behaviour change require separate studies.

\section{Conclusion}
Prompt design and early stopping explain much of SgFood-Edit's apparent gain. Retry passes the checker on 28/34 cases, versus 26/34 for best-of-four and fixed-prompt stopping; the latter matches retry's 68 generations. The task-dependent difference is uncertain. CLIP and the fixed recolouring control broaden diagnostics, not quality evidence. Human ratings, checker-error analysis and new-photo confirmation remain necessary.

{\small
\bibliographystyle{ieeenat_fullname}
\bibliography{main}
}
\end{document}
